\subsubsection{Production Considerations}

\begin{principle} \hlt{Trading Architecture}\\
Automated trading needs a production architecture that synchronises market data, order gateways, historical databases, models, monitoring, and TCA. The hard part is not one model; it is keeping the decision chain coherent under fragmented venues and live latency.
\end{principle}

\begin{taxonomy} \hlt{Trading Tool Input Segmentation}
\begin{enumerate}[label=\roman*.]
\setlength{\itemsep}{0pt}
\item \hlt{market data}: trades, quotes, depth, timestamps, and book updates, preferably consolidated across venues;
\item \hlt{venue connectivity}: private order-entry messages and public market-data feeds;
\item \hlt{historical data}: raw events, real-time analytics, order messages, fills, and market context;
\item \hlt{models}: impact, risk, alpha, volume, and context models with calibrated parameters;
\item \hlt{calendar}: holidays, auctions, expiries, index events, and expected halt days.
\end{enumerate}
\end{taxonomy}

\begin{taxonomy} \hlt{Broker-Dealer Operations Segmentation}
\begin{enumerate}[label=\roman*.]
\setlength{\itemsep}{0pt}
\item \hlt{front office}: client contact, order receipt, execution, sales, and research;
\item \hlt{back office}: accounts, confirmations, clearing, settlement, custody, and statements;
\item \hlt{risk and treasury}: margin, credit, stock loan, cash management, and firm inventory;
\item \hlt{compliance}: suitability, best execution, surveillance, and regulatory reporting;
\item \hlt{technology}: order routing, market data, drop copies, books, and audit trails.
\end{enumerate}
\end{taxonomy}

\begin{mechanism} \hlt{Analytics Close to the Feed}\\
Spread, effective spread, volatility, book imbalance, and queue-size measures should be computed as soon as raw data arrive. Fast analytics become state variables for execution, monitoring, and later post-trade explanation.
\end{mechanism}

\begin{caveat} \hlt{Connections Are Two-Way}\\
An execution engine sends orders, cancellations, and modifications, then receives acknowledgements, rejects, fills, and public book updates. In dark venues, the private responses can also act as probes of hidden liquidity.
\end{caveat}

\begin{taxonomy} \hlt{Model Component Segmentation}
\begin{enumerate}[label=\roman*.]
\setlength{\itemsep}{0pt}
\item \hlt{formula}: the functional relation between state variables and decisions;
\item \hlt{estimation}: the procedure that calibrates parameters from historical data;
\item \hlt{accuracy check}: the live or post-trade test that decides whether the model is still usable.
\end{enumerate}
\end{taxonomy}

\begin{process} \hlt{Indicator Publication}
\begin{enumerate}[label=\roman*.]
\setlength{\itemsep}{0pt}
\item models and algorithms publish their state as indicators;
\item traders monitor exceptions, concentration, performance, and context shifts;
\item post-trade systems store the same indicators for diagnosis and TCA.
\end{enumerate}
This makes trading algorithms stateful agents in the production stack rather than black boxes.
\end{process}

\begin{process} \hlt{Audit Trail}
\begin{enumerate}[label=\roman*.]
\setlength{\itemsep}{0pt}
\item record order receipt, route, modification, cancellation, execution, allocation, and timestamp;
\item link private order messages to public market state and venue responses;
\item preserve enough context to test best execution, front-running, allocation, and unauthorised trading.
\end{enumerate}
\end{process}

\begin{taxonomy} \hlt{Automated Trading Reliability Segmentation}
\begin{enumerate}[label=\roman*.]
\setlength{\itemsep}{0pt}
\item \hlt{availability}: redundant servers, gateways, power, and network paths;
\item \hlt{integrity}: authenticated users, validated orders, and tamper-resistant logs;
\item \hlt{recoverability}: backups, replayable feeds, drop copies, and disaster-recovery sites;
\item \hlt{capacity}: headroom for message bursts, opens, closes, halts, and data-feed spikes;
\item \hlt{latency}: bounded delays for market data, decisions, order entry, and acknowledgements.
\end{enumerate}
For a trading system, reliability is a market-quality variable, not only an IT variable.
\end{taxonomy}

\begin{caveat} \hlt{Network Failure Creates Position Risk}\\
If order entry, market data, or drop copy fails, the trader may not know whether an order traded, whether a hedge executed, or whether the market moved through a stop.\\
Controls should prefer safe state reconciliation over blind resubmission.
\end{caveat}

\begin{mechanism} \hlt{Automated System Cost Structure}\\
Electronic trading has high fixed costs for software, data networks, security, and recovery, but low marginal labour cost per order. Floor trading has higher human and physical operating costs, but can still be valuable when negotiation, discretion, and soft information dominate raw matching speed.
\end{mechanism}
